\documentclass[11pt]{article}
\usepackage{mystyle}

\title{Rosen --- \S 6.4: Binomial Coefficients and Identities\\ \large Selected Exercises}
\author{Alston}
\date{Semester 2, 2026}

\begin{document}
\maketitle

% ======================================================================
\begin{exercise}{21}
    Prove that if $n$ and $k$ are integers with $1 \leq k \leq n$,
    then $k\binom{n}{k} = n\binom{n-1}{k-1}$,
    \begin{enumerate}[label=(\alph*)]
        \item using a combinatorial proof.
        [\emph{Hint}: Show that the two sides of the identity count
        the number of ways to select a subset with $k$ elements from
        a set with $n$ elements and then an element of this subset.]
        \item using an algebraic proof based on the formula for
        $\binom{n}{r}$ given in Theorem 2 in Section 6.3.
    \end{enumerate}
\end{exercise}

% ======================================================================
\begin{exercise}{24}
    Show that if $p$ is a prime and $k$ is an integer such that
    $1 \leq k \leq p-1$, then $p$ divides $\binom{p}{k}$.
\end{exercise}

% ======================================================================
\begin{exercise}{27}
    Prove the \emph{hockeystick identity}
    \[
    \sum_{k=0}^{r} \binom{n+k}{k} = \binom{n+r+1}{r}
    \]
    whenever $n$ and $r$ are positive integers,
    \begin{enumerate}[label=(\alph*)]
        \item using a combinatorial argument.
        \item using Pascal's identity.
    \end{enumerate}
\end{exercise}

% ======================================================================
\begin{exercise}{31}
    Show that a nonempty set has the same number of subsets with an
    odd number of elements as it does subsets with an even number of
    elements.
\end{exercise}

% ======================================================================
\begin{exercise}{32}
    Prove the binomial theorem using mathematical induction.
\end{exercise}

% ======================================================================
\begin{exercise}{33}
    In this exercise we will count the number of paths in the $xy$
    plane between the origin $(0,0)$ and point $(m,n)$, where $m$
    and $n$ are nonnegative integers, such that each path is made up
    of a series of steps, where each step is a move one unit to the
    right or a move one unit upward. (No moves to the left or
    downward are allowed.)
    \begin{enumerate}[label=(\alph*)]
        \item Show that each path of the type described can be
        represented by a bit string consisting of $m$ 0s and $n$ 1s,
        where a 0 represents a move one unit to the right and a 1
        represents a move one unit upward.
        \item Conclude from part (a) that there are $\binom{m+n}{n}$
        paths of the desired type.
    \end{enumerate}
\end{exercise}

% ======================================================================
\begin{exercise}{36}
    Use Exercise 33 to prove Pascal's identity.
    [\emph{Hint}: Show that a path of the type described in Exercise
    33 from $(0,0)$ to $(n+1-k, k)$ passes through either
    $(n+1-k, k-1)$ or $(n-k, k)$, but not through both.]
\end{exercise}

% ======================================================================
\begin{exercise}{38}
    Give a combinatorial proof that if $n$ is a positive integer then
    \[
    \sum_{k=0}^{n} k^2 \binom{n}{k} = n(n+1)2^{n-2}.
    \]
    [\emph{Hint}: Show that both sides count the ways to select a
    subset of a set of $n$ elements together with two not necessarily
    distinct elements from this subset. Furthermore, express the
    right-hand side as $n(n-1)2^{n-2} + n2^{n-1}$.]
\end{exercise}

\end{document}